%% ============================================================
%% shared/architecture_patterns.tex
%% Key AWS Architecture Patterns
%% Classification: BOTH (SAA: know all; SAP: design and justify each)
%% ============================================================

\servicesection{Key Architecture Patterns}{\bothtag}{%
  Reference blueprints for the most common exam scenario architectures.}

These patterns appear repeatedly across both exams. For each one, know the
services involved, why each service is chosen, and what the exam alternatives
look like when constraints change (e.g.\ lowest cost, lowest RTO, no containers).

%% ─────────────────────────────────────────────────────────────
\subsection{Three-Tier Web Application}
%% ─────────────────────────────────────────────────────────────

\begin{keyfacts}
\begin{itemize}
  \item \textbf{Presentation tier}: Route 53 (DNS) $\to$ CloudFront (CDN/cache) $\to$ ALB
        (Layer-7 routing, SSL termination).
  \item \textbf{Application tier}: Auto Scaling Group of EC2 in \emph{private subnets}
        across $\geq$2 AZs; or ECS Fargate tasks (no cluster management).
  \item \textbf{Data tier}: RDS Multi-AZ (primary write) + Read Replicas (read scaling) +
        ElastiCache Redis (session/object cache to reduce DB load).
  \item Static assets (HTML/CSS/JS/images) served from \textbf{S3 via CloudFront with OAC}
        — eliminates direct S3 public access.
  \item Use \textbf{S3 Gateway Endpoint} so EC2 instances reach S3 without traversing the
        internet (no NAT Gateway cost, no data-transfer charge).
\end{itemize}
\end{keyfacts}

\begin{examtip}[title={[!] When the question changes the constraint}]
\begin{itemize}
  \item \textbf{``Reduce DB read load''} $\to$ add ElastiCache Redis (cache-aside) or
        RDS Read Replica.
  \item \textbf{``Higher availability across regions''} $\to$ add Route 53 Failover + warm
        standby in second region; upgrade DB to Aurora Global Database.
  \item \textbf{``Least operational overhead''} $\to$ replace EC2 ASG with ECS Fargate;
        replace RDS with Aurora Serverless v2.
  \item \textbf{``Static IP required''} $\to$ replace ALB with NLB; add Elastic IP.
\end{itemize}
\end{examtip}

%% ─────────────────────────────────────────────────────────────
\subsection{Serverless Web Application}
%% ─────────────────────────────────────────────────────────────

\begin{keyfacts}
\begin{itemize}
  \item \textbf{Full stack}: Route 53 $\to$ CloudFront $\to$ API Gateway $\to$ Lambda $\to$
        DynamoDB (+ DAX for sub-ms reads).
  \item \textbf{Static content}: S3 bucket serving the SPA/frontend via CloudFront with OAC.
  \item \textbf{Auth}: Amazon Cognito User Pools (user directory + JWT tokens) + Cognito
        Identity Pools (exchange token for temporary AWS credentials).
  \item \textbf{No servers to patch}; scales automatically to zero and to millions; pay per
        request only.
  \item Lambda@Edge or CloudFront Functions can handle auth, redirects, and header
        manipulation at the edge — before the request reaches API Gateway.
\end{itemize}
\end{keyfacts}

\begin{comparebox}[title={Serverless web app --- service substitutions}]
\begin{tabular}{@{}p{4.5cm}p{7cm}@{}}
\toprule
\textbf{Need} & \textbf{Service} \\
\midrule
HTTP API, low latency, cheap          & API Gateway HTTP API (not REST API; ~70\% cheaper) \\
REST API with full features (throttle, cache, keys) & API Gateway REST API \\
Simple web app / API, full managed    & AWS App Runner (from container or source) \\
Background async processing           & Lambda triggered by SQS / EventBridge \\
Long-running tasks ($>$15 min)        & ECS Fargate task (replace Lambda) \\
Complex multi-step workflow           & AWS Step Functions \\
\bottomrule
\end{tabular}
\end{comparebox}

%% ─────────────────────────────────────────────────────────────
\subsection{Event-Driven / Decoupled Pipeline}
%% ─────────────────────────────────────────────────────────────

\begin{keyfacts}
\begin{itemize}
  \item \textbf{Pattern}: Producer $\to$ SQS / Kinesis Data Streams $\to$ Lambda / ECS
        workers $\to$ S3 / DynamoDB.
  \item \textbf{Fan-out}: SNS topic $\to$ multiple SQS queues; each queue drives an
        independent downstream Lambda or worker — allows parallel processing without
        coordination.
  \item \textbf{Rule-based routing}: EventBridge event bus; pattern rules filter events to
        specific Lambda/SQS/Step Functions targets; preferred over SNS for event-driven
        architectures reacting to AWS service events or SaaS.
  \item \textbf{Failure handling}: Dead-letter queue (DLQ) on SQS or Lambda; messages that
        exceed \texttt{maxReceiveCount} land in DLQ for inspection.
  \item \textbf{Ordered + replay}: Kinesis Data Streams (up to 7-day replay, ordered per
        shard); use when consumers need to replay or process in order.
\end{itemize}
\end{keyfacts}

\begin{examtip}[title={[!] SQS vs Kinesis decision}]
\begin{itemize}
  \item \textbf{SQS}: decouple, buffer, unknown order OK, each message processed once $\to$
        \textbf{SQS Standard}.
  \item \textbf{SQS FIFO}: strict order, exactly-once processing (financial transactions,
        inventory deduction).
  \item \textbf{Kinesis Data Streams}: ordered per shard, replay up to 7 days, multiple
        consumers reading concurrently (fan-out without copying messages).
  \item \textbf{Kinesis Firehose}: fully managed delivery to S3/Redshift/OpenSearch; no
        consumer code needed; near-real-time (60 s buffer or 1 MB).
\end{itemize}
\end{examtip}

%% ─────────────────────────────────────────────────────────────
\subsection{Data Lake Architecture}
%% ─────────────────────────────────────────────────────────────

\begin{keyfacts}
\begin{itemize}
  \item \textbf{Ingestion}: Kinesis Data Firehose (streaming) or AWS DataSync (bulk online)
        or Snow Family (offline) $\to$ S3 raw/landing zone.
  \item \textbf{Catalogue \& ETL}: AWS Glue (serverless ETL + Data Catalog) transforms raw
        data $\to$ S3 curated/processed zone in Parquet/ORC.
  \item \textbf{Query}: Amazon Athena (SQL on S3; pay per query; no infrastructure) or
        Redshift Spectrum (query S3 directly from a Redshift cluster).
  \item \textbf{Visualise}: Amazon QuickSight (BI dashboards; connects to Athena/Redshift/S3).
  \item \textbf{Access control}: S3 bucket policies + Lake Formation (column/row-level
        security on the data catalog).
\end{itemize}
\end{keyfacts}

\begin{examtip}[title={[!] Data lake query pattern shortcuts}]
\begin{itemize}
  \item \textbf{``Ad-hoc queries on S3 without loading''} $\to$ \textbf{Athena}.
  \item \textbf{``Join S3 data with Redshift warehouse''} $\to$ \textbf{Redshift Spectrum}.
  \item \textbf{``Centralise metadata and permissions across all data''} $\to$
        \textbf{AWS Lake Formation}.
  \item \textbf{``Serverless ETL with dependency tracking''} $\to$ \textbf{AWS Glue}.
\end{itemize}
\end{examtip}

%% ─────────────────────────────────────────────────────────────
\subsection{Hybrid Cloud Pattern}
%% ─────────────────────────────────────────────────────────────

\begin{keyfacts}
\begin{itemize}
  \item \textbf{Network}: on-prem $\leftrightarrow$ Direct Connect (consistent, private,
        lower latency) \emph{or} Site-to-Site VPN (lower cost, internet, encrypted).
        Use \textbf{both} for redundancy (VPN as failover for DX).
  \item \textbf{File storage extension}: Storage Gateway (File mode) $\to$ S3; local cache
        keeps hot files fast; cloud stores everything.
  \item \textbf{Automated data transfer}: AWS DataSync (scheduled/incremental; NFS/SMB/HDFS
        $\to$ S3/EFS/FSx); up to 10 Gbps per task.
  \item \textbf{Bulk offline transfer}: Snowball Edge (terabytes; 80 TB per device) or
        Snowmobile (petabytes; 100 PB per truck) when bandwidth is insufficient.
  \item \textbf{DNS}: Route 53 Resolver inbound/outbound endpoints allow private DNS
        resolution between on-prem DNS and Route 53 Private Hosted Zones.
  \item \textbf{Identity}: AWS Directory Service (Managed Microsoft AD) for SSO between
        on-prem Active Directory and AWS resources.
\end{itemize}
\end{keyfacts}

\begin{comparebox}[title={Direct Connect vs VPN}]
\begin{tabular}{@{}p{3.5cm}p{4.5cm}p{3.5cm}@{}}
\toprule
 & \textbf{Direct Connect} & \textbf{Site-to-Site VPN} \\
\midrule
Path       & Dedicated private fibre              & Encrypted over internet \\
Latency    & Consistent, low                      & Variable \\
Bandwidth  & 1 Gbps -- 100 Gbps                   & Up to 1.25 Gbps \\
Cost       & Higher (port + data)                 & Lower \\
Setup time & Weeks                                & Minutes \\
Encryption & Not by default (add VPN on DX)       & Always encrypted (IPsec) \\
Best for   & Steady high-volume, compliance       & Backup/failover, quick setup \\
\bottomrule
\end{tabular}
\end{comparebox}

%% ─────────────────────────────────────────────────────────────
\subsection{Migration Services Reference}
%% ─────────────────────────────────────────────────────────────

\begin{comparebox}[title={Migration tool selection}]
\begin{tabular}{@{}p{4cm}p{7.5cm}@{}}
\toprule
\textbf{Service} & \textbf{Use case} \\
\midrule
AWS DMS (Database Migration Service) & Migrate relational/NoSQL DBs; heterogeneous migrations; CDC for minimal downtime \\[3pt]
AWS SCT (Schema Conversion Tool)     & Convert schema from one engine to another (Oracle $\to$ Aurora); runs alongside DMS \\[3pt]
AWS DataSync                         & Online automated transfer; NFS/SMB/HDFS $\to$ S3/EFS/FSx; up to 10 Gbps \\[3pt]
AWS Snowball Edge                    & Offline bulk transfer; 80 TB usable storage; compute at edge option \\[3pt]
AWS Snowmobile                       & Exabyte-scale ($\geq$100 PB); physical truck transport \\[3pt]
AWS Application Migration Service (MGN) & Lift-and-shift; continuous block-level replication; minimal-downtime cutover \\[3pt]
AWS Transfer Family                  & Managed SFTP/FTPS/FTP/AS2 to S3 or EFS; drop-in replacement for on-prem FTP servers \\[3pt]
\bottomrule
\end{tabular}
\end{comparebox}

\begin{examtip}[title={[!] Migration exam shortcuts}]
\begin{itemize}
  \item \textbf{``Transfer 2 PB with limited bandwidth''} $\to$ \textbf{Snowball Edge} (multiple devices).
  \item \textbf{``Transfer 100 PB or more''} $\to$ \textbf{Snowmobile}.
  \item \textbf{``Ongoing DB migration, minimal downtime''} $\to$ \textbf{DMS with CDC} (Change Data Capture).
  \item \textbf{``Oracle to Aurora''} $\to$ \textbf{SCT} (schema) + \textbf{DMS} (data).
  \item \textbf{``Lift-and-shift 200 servers''} $\to$ \textbf{Application Migration Service (MGN)}.
  \item \textbf{``Automated periodic transfer from on-prem NAS to S3''} $\to$ \textbf{DataSync}.
\end{itemize}
\end{examtip}

\begin{sapdepth}
\textbf{SAP migration depth}: SAP-C02 adds the \textbf{7Rs framework} (Retire, Retain, Rehost,
Relocate, Replatform, Repurchase, Refactor) as a \emph{decision methodology}, not just a tool list.
Refactor (re-architect to cloud-native) has the highest long-term benefit but the highest effort.
Rehost (lift-and-shift via MGN) has the lowest benefit but lowest effort and is the fastest path.
SAP scenarios will ask you to recommend the right ``R'' for a given business constraint.
\end{sapdepth}
